\scanpage{649}\mspage{37}
$f=pi$ \uline{fonctorielle} en $f$, ce qui est le cas
en particulier si la structure à cofibrations
est spéciale. Alors pour \unsure{tt|tte} $I$ dans
Cat, $\deleted{\unsure{M}}M(I)$ muni de $(W(I),C(I))$ est
une catégorie à cofibrations, laquelle \deleted{\unsure{ne peut}}
satisfait l'axiome transfini Cof 4 si
$(W,C)$ y satisfait. Donc on peut lui
appliquer les résultats de 3.2, qui
impliquent notamment que les axiomes
Der $1,2,3\,(?)$ sont valables sur Cat \unsure{tout entier|toute entière}, et que $u^{W}_{!}$ existe quand $u$
est de la forme $\mathrm{id}_{I}\times v$, avec $v:J\unsure{-|\to}K$
une flèche \unsure{dans} $\mathrm{Diag}_1$ (resp. dans $\mathrm{Diag}_2$, si
l'axiome transfini est satisfait).

3.4. La clef de la théorie est dans l'extension
de la structure à \unsure{cofibration|cofibrations} $(W,C)$, de $M$
à une catégorie $M(I)$, où $I$ est un ensemble
ordonné, en $(W(I),C_I)$. Les faits principaux
sont énoncés dans \inserted{$\underline{\overline{\mathrm{XVII}}}$, \unsure{aux|en}} \deleted{\unsure{les}} théorèmes p.~26--27, p.~50, p.~60.
C'est quand $I$ est \uline{noethérien}, et qu'on suppose
soit que $\mathrm{Esp}(I^\circ)$ est également noethérien, (p.ex.
$I$ fini), soit que $(W,C)$ satisfait l'axiome
transfini Cof 4, qu'on prouve que $(W(I),C_I)$
est également une structure à cofibrations sur $M(I)$,
% archive sheet number: 648

